\documentclass[12pt]{article}

% revise margins
\setlength{\headheight}{0.0in}
\setlength{\topmargin}{0.0in}
\setlength{\headsep}{0.0in}
\setlength{\textheight}{9.0in}
\setlength{\footskip}{0.0in}
\setlength{\oddsidemargin}{0.0in}
\setlength{\evensidemargin}{0.0in}
\setlength{\textwidth}{6.5in}

\usepackage{times}
\usepackage{amsmath}

\begin{document}

\thispagestyle{empty}

\noindent \textbf{Biostat 140.668} \hfill Solutions to Problem set 1

\bigskip 

\noindent \textbf{Meiosis, recombination, interference}

\begin{enumerate}

\item Let $m$ denote the number of crossovers in an interval of length $d$
Morgans.  Then $m \sim \text{Poisson}(d)$.  The map function expresses
the probability of a recombination event as a function of the genetic
distance $d$.  But the probability of a recombination event is simply
the probability that $m$ is odd.

$$\Pr(m \text{ is odd})  =  \sum_{k=0}^{\infty} \Pr(m = 2k+1) 
 =  \sum_{k=0}^{\infty} e^{-d} \frac{d^{2k+1}}{(2k+1)!}$$

Note that 
\begin{eqnarray*}
e^d = \sum_{k=0}^{\infty} d^k/k! & = & 1 + d + d^2/(2!) +
d^3/(3!) + \dots \\
e^{-d} = \sum_{k=0}^{\infty} (-1)^k d^k/k! & = & 1
- d + d^2/(2!) - d^3/(3!) + \dots
\end{eqnarray*}

Thus $e^d - e^{-d} = 2(d + d^3/(3!) + d^5/(5!) + \dots) = 2
\sum_{k=0}^{\infty} d^{2k+1}/(2k+1)!$.

And so $r = \Pr(m \text{ is odd}) = e^{-d} (e^{d} - e^{-d})/2 =
(1-e^{-2d})/2$, which is the Haldane map function.
\bigskip

\item This is rather a cute problem.  If you toss $n \ge 1$ fair coins, the
chance of obtaining an odd number of heads is 1/2.  

The simplest solution is to use induction.  

The cases $n=0$ and $n=1$ are obvious.  

Suppose you toss $n+1$ fair coins, where $n \ge 1$.  Pr(odd number of
heads in $n+1$ tosses) = Pr(odd number of heads in $n$ tosses)
$\times$ Pr($n+1$st toss is tails) + Pr(even number of heads in $n$
tosses) $\times$ Pr($n+1$st toss is heads) = \ldots = 1/2.


\item This is a bit tricky.  

Assume that the first chiamsa involves a random pair of non-sister
chromatids, that the next chiasma involves exactly the opposite pair,
and subsequent chiasmata alternate in their choice of strands.  This
is strong chromatid interference.

Assume that the locations of chiasmata on the four-strand bundle are
according to a stationary renewal process with inter-arrival
distribution gamma(shape=1/2, rate=1).  This is strong negative
chiasma interference.

Recall that if $X_1, X_2 \sim$ iid gamma(shape=1/2, rate=1), then $X_1
+ X_2 \sim$ gamma(shape=1, rate=1) $\equiv$ exponential(rate=1).  

It should be clear, then, that the locations of crossovers on a random
meiotic product follow a Poisson process, and so exhibit no crossover
interference.


\end{enumerate}




\end{document}
